\documentclass[10pt]{article}
\usepackage[T1]{fontenc}
\usepackage[utf8]{inputenc}
\usepackage[french]{babel}
\usepackage{amsmath,amssymb,mathrsfs,mathpazo}
\usepackage{geometry}
\geometry{a4paper,margin=2cm}
\usepackage{graphicx,array,tabularx,tabularray,colortbl,multirow,multicol,booktabs,enumitem,pifont,ulem,fancybox,diagbox,cancel}
\usepackage{pstricks,pstricks-add,pst-plot,pst-tree,pst-node,pst-text,pst-func,pst-math,pst-bspline,pst-3dplot,pst-eucl,pst-solides3d}
\let\pscalculate\undefined
\usepackage{pst-bezier}
\usepackage{tikz,pgfplots,tkz-tab}
\usetikzlibrary{arrows,arrows.meta,patterns,calc,positioning,trees,decorations.markings,shapes.geometric,shapes.misc,shapes.arrows,intersections,decorations.pathmorphing,decorations.pathreplacing,backgrounds,shadows,babel}
\pgfplotsset{compat=1.18}
\usepgfplotslibrary{fillbetween}
\usepackage[np]{numprint}
\usepackage{listings,makecell,pas-tableur,twoopt}
\let\barre\undefined
\usepackage[pstricks]{tablvar}
\definecolor{darkgreen}{RGB}{0,100,0}
\providecommand{\red}{\color{red}}
\providecommand{\blue}{\color{blue}}
\providecommand{\green}{\color{green}}
\providecommand{\black}{\color{black}}
\providecommand{\white}{\color{white}}
\providecommand{\gray}{\color{gray}}
\providecommand{\cyan}{\color{cyan}}
\providecommand{\magenta}{\color{magenta}}
\providecommand{\yellow}{\color{yellow}}
\providecommand{\violet}{\color{violet}}
\providecommand{\orange}{\color{orange}}
\providecommand{\euro}{\texteuro}
\providecommand{\EUR}{\texteuro}
\providecommand{\ieme}{\textsuperscript{e}}
\providecommand{\degres}{\ensuremath{^\circ}}
\providecommand{\degree}{\ensuremath{^\circ}}
\providecommand{\vv}[1]{\overrightarrow{#1}}
\providecommand{\Coord}[1]{#1}
\providecommand{\decosix}{}
\providecommand{\starredbullet}{\ensuremath{\bullet}}
\pagestyle{empty}
\setlength{\parindent}{0pt}
\usepackage{adjustbox}\DeclareUnicodeCharacter{2236}{:}\DeclareUnicodeCharacter{2190}{\ensuremath{\leftarrow}}
\providecommand{\R}{\mathbb{R}}
\providecommand{\N}{\mathbb{N}}
\providecommand{\Z}{\mathbb{Z}}
\providecommand{\Q}{\mathbb{Q}}
\providecommand{\D}{\mathbb{D}}
\providecommand{\C}{\mathbb{C}}
\providecommand{\e}{{\mathrm{e}}}
\providecommand{\vect}[1]{\overrightarrow{#1}}
\providecommand{\vectt}[1]{\overrightarrow{\mathrm{#1}}}
\providecommand{\vvt}[1]{\overrightarrow{\mathrm{#1}}}
\providecommand{\np}[2][]{{#2~#1}}
\providecommand{\vv}[1]{{\overrightarrow{#1}}}
\renewcommand{\d}{{\mathrm{d}}}
\providecommand{\degres}{{^\circ}}
\providecommand{\degre}{{^\circ}}
\providecommand{\blue}{\color{blue}}
\providecommand{\red}{\color{red}}
\providecommand{\enskip}{\quad}
\providecommand{\llbracket}{\lbrack\!\lbrack}
\providecommand{\rrbracket}{\rbrack\!\rbrack}
\providecommand{\euro}{{\text{€}}}
\providecommand{\no}{{\text{n°}}}
\providecommand{\interval}[2]{{[#1;#2]}}
\providecommand{\intervFO}[2]{{[#1;#2[}}
\providecommand{\dotfill}{\ldots}

\providecommand{\ds}{\displaystyle}
\providecommand{\textding}[1]{#1}
\providecommand{\Oij}{\(\mathrm{O};\vec{\imath},\vec{\jmath}\)}
\providecommand{\Oijk}{\(\mathrm{O};\vec{\imath},\vec{\jmath},\vec{k}\)}
\providecommand{\Ouv}{\(\mathrm{O};\vec{u},\vec{v}\)}
\providecommand{\intervFF}[2]{\ensuremath{[#1;#2]}}
\renewcommand{\euro}{\ifmmode\text{\texteuro}\else\texteuro\fi}
\DeclareUnicodeCharacter{00B0}{\ifmmode^{\circ}\else\textdegree\fi}
\DeclareUnicodeCharacter{20AC}{\ifmmode\text{\texteuro}\else\texteuro\fi}
% Énoncés : Éducation nationale. Transcription et corrigés : APMEP, Denis Vergès et contributeurs. Découpage et classement : CoopMaths, Sébastien Lozano, Éric Elter et contributeurs. Adaptation de la présentation : Maths Explorer.
% CC BY-SA : https://coopmaths.fr/www/about/
% https://coopmaths.fr/alea/static/bac/2021/tex/bac_2021_03_sujet1_3_cor.tex
\begin{document}
\noindent\textit{APMEP / CoopMaths. CC BY-SA. Présentation adaptée par Maths Explorer.}\par\medskip
\begin{center}
\psset{unit=1cm}
\begin{pspicture}(-0.5,-0.5)(9,5)
\psline(0,0)(5,0)(8,1.8)
\pspolygon[linestyle=dashed](0,0)(8,1.8)(3,1.8)(0,0)
\psline[linestyle=dashed](5,0)(3,1.8)(4,4)(4,0.95)
\psline(0,0)(4,4)(5,0)(8,1.8)(4,4)
\uput[dl](0,0){D}\uput[dr](5,0){C}\uput[ur](8,1.8){B}\uput[ul](3,1.8){A}
\uput[u](4,4){S}\uput[d](4,0.95){I}
\uput[ul](2,2){K}\uput[ur](4.5,2){L}\uput[ur](6,2.9){M}
\psdots[dotstyle=+,dotangle=30,dotscale=1.85](2,2)(4.5,2)(6,2.9)
\end{pspicture}
\end{center}

SABCD est une pyramide régulière à base carrée ABCD dont toutes les arêtes ont la même longueur.
Le point I est le centre du carré ABCD. 
On suppose que : IC = IB = IS $= 1$.

Les points K, L et M sont les milieux respectifs des arêtes [SD], [SC] et [SB].



\begin{enumerate}
\item Les droites suivantes ne sont pas coplanaires:

\begin{center}
\begin{tabularx}{\linewidth}{*{4}{X}}
\textbf{a.~} (DK) et (SD) &\textbf{b.~} (AS) et (IC) &\fbox{\textbf{c.~} (AC) et (SB)} &\textbf{d.~} (LM) et (AD)
\end{tabularx}
\end{center}

\begin{tabular}{@{\hspace*{0.05\linewidth}} | p{0.93\linewidth}}
\vspace*{-10pt}
\begin{list}{\textbullet}{On procède par élimination.}
\item Les droites (DK) et (SD) sont sécantes en D donc  coplanaires; on élimine \textbf{a.}
\item Les droites (AS) et (IC) sont sécantes en A donc coplanaires; on élimine \textbf{b.}
\item Les droites (LM) et (AD) sont toutes deux parallèles à (BC) donc parallèles entre elles; elles sont donc coplanaires; on élimine \textbf{d.}
\end{list}

\smallskip

\textbf{Réponse c.}
\end{tabular}
\end{enumerate}

Pour les questions suivantes, on se place dans le repère orthonormé de l'espace $\left(\text{ I}~;~ \vect{\text{IC}},~\vect{\text{IB}},~\vect{\text{IS}}\right)$.

Dans ce repère, on donne les coordonnées des points suivants:

\[\text{I}(0~;~0~;~0) \:;\: \text{A}(-1~;~0~;~0) \:;\: \text{B}(0~;~1~;~0) \:;\: \text{C}(1~;~0~;~0); \text{D}(0~;~-1~;~0) \:;\: \text{S}(0~;~0~;~1).\]

\begin{enumerate}[resume]
\item  Les coordonnées du milieu N de [KL] sont:
\begin{center}
\begin{tabularx}{\linewidth}{*{4}{X}}
\textbf{a.~} $\left(\dfrac{1}{4}~;~\dfrac{1}{4}~;~\dfrac{1}{4}\right)$&
\fbox{\textbf{b.~}$\left(\dfrac{1}{4}~;~- \dfrac{1}{4}~;~\dfrac{1}{2}\right)$}&
\textbf{c.~}$\left(-\dfrac{1}{4}~;~\dfrac{1}{4}~;~\dfrac{1}{2}\right)$&
\textbf{d.~}$\left(-\dfrac{1}{2}~;~\dfrac{1}{2}~;~1\right)$
\end{tabularx}
\end{center}

\begin{tabular}{@{\hspace*{0.05\linewidth}} | p{0.93\linewidth}}
\vspace*{-10pt}
\begin{list}{\textbullet}{}
\item Le milieu K de [SD] a pour coordonnées $\left ( 0~;~-\frac{1}{2}~;~\frac{1}{2}\right )$.
\item Le milieu L de [SC] a pour coordonnées $\left ( \frac{1}{2}~;~0~;~\frac{1}{2}\right )$.
\item Le milieu N de [KL] a donc pour coordonnées $\left ( \frac{1}{4}~;~-\frac{1}{4}~;~\frac{1}{2}\right )$.
\end{list}

\smallskip

\textbf{Réponse b.}
\end{tabular}

\medskip

\item  Les coordonnées du vecteur $\vect{\text{AS}}$ sont:
\begin{center}
\begin{tabularx}{\linewidth}{*{4}{X}}
\textbf{a.~} $\begin{pmatrix}1\\1\\0 \end{pmatrix}$&
\fbox{\textbf{b.~}  $\begin{pmatrix}1\\0\\1 \end{pmatrix}$}&
\textbf{c.~} $\begin{pmatrix}2\\1\\-1 \end{pmatrix}$ &
\textbf{d.~} $\begin{pmatrix} 1\\1\\1\end{pmatrix}$
\end{tabularx}
\end{center}

\begin{tabular}{@{\hspace*{0.05\linewidth}} | p{0.93\linewidth}}
\textbf{Réponse b.}
\end{tabular}

\medskip

\item Une représentation paramétrique de la droite (AS) est:

{\footnotesize \begin{center}
\begin{tabularx}{\linewidth}{*{4}{X}}
\textbf{a.~} $\left\{\begin{array}{l c r}x&=&-1-t\\y&=&t\\z&=&-t
\end{array}\right.$

$(t \in \R)$&\textbf{b.~} $\left\{\begin{array}{l c r}x&=&-1+2t\\y&=&0\\z&=&1 + 2t
\end{array}\right.$

$(t \in \R)$&
\fbox{\parbox{2.6cm}{\textbf{c.~} $\left\{\begin{array}{l c r}x&=&t\\y&=&0\\z&=&1+t
\end{array}\right.$\\
$(t \in \R)$}}
&\textbf{d.~} $\left\{\begin{array}{l c r}x&=&-1-t\\y&=&1+t\\z&=&1-t
\end{array}\right.$

$(t \in \R)$
\end{tabularx}
\end{center}}

\begin{tabular}{@{\hspace*{0.05\linewidth}} | p{0.93\linewidth}}
La droite (AS) a pour vecteur directeur $\vectt{AS}\,(1~;~0~;~1)$; la seule représentation qui convienne est la c.

\smallskip

\textbf{Réponse c.}
\end{tabular}

 \medskip
 
\item Une équation cartésienne du plan (SCB) est:

\begin{center}
\begin{tabularx}{\linewidth}{*{4}{X}}
\textbf{a.~} $y+z-1 =0$ &\fbox{\textbf{b.~}$x+y+z- 1=0$} & \textbf{c.~}$x-y+z=0$&
\textbf{d.~}$x+z-1 =0$
\end{tabularx}
\end{center}

\begin{tabular}{@{\hspace*{0.05\linewidth}} | p{0.93\linewidth}}
\vspace*{-10pt}
\begin{list}{\textbullet}{On procède par élimination.}
\item Le plan d'équation $y+z-1=0$ ne contient pas  C\,$(1~;~0~;~0)$; on élimine \textbf{a.}
\item Le plan d'équation $x-y+z=0$ ne contient pas  S\,$(0~;~0~;~1)$; on élimine \textbf{c.}
\item Le plan d'équation $x+z-1=0$ ne contient pas  B\,$(0~;~1~;~0)$; on élimine \textbf{d.}
\end{list}

\smallskip

\textbf{Réponse b.}
\end{tabular}

\end{enumerate}

\bigskip

\end{document}
